Los resultados muestran que el desempeño práctico depende tanto del tamaño como
de la distribución de los datos. En sorting, \texttt{std::sort} fue el algoritmo más
rápido en los tamaños pequeños y medianos, mientras que QuickSort y
PatienceSort fueron particularmente sensibles a entradas con muchos valores
repetidos o a ciertos órdenes iniciales. En el caso $10^7$ medido, PatienceSort
alcanzó aproximadamente 380 segundos, frente a 0.68 segundos de \texttt{std::sort}.
En matrices, Strassen fue más lento que Naive en las dimensiones probadas,
aunque ambos entregaron resultados iguales; esto muestra que una mejor
complejidad asintótica no garantiza una ventaja para tamaños donde dominan las
constantes y la gestión de memoria.
